\documentclass[10pt,a4paper]{amsart}
\usepackage{amsmath,amssymb,amsthm,graphicx}
\usepackage{verbatim}
\newenvironment{reference}{\comment}{\endcomment}
%\newenvironment{reference}{}{}
\newenvironment{slogan}{\comment}{\endcomment}
\newenvironment{history}{\comment}{\endcomment}

% For commutative diagrams we use Xy-pic
\usepackage[all]{xy}

% We use 2cell for 2-commutative diagrams.
\xyoption{2cell}
% Theorem environments.
%
\theoremstyle{plain}
\newtheorem{theorem}[subsection]{Theorem}
\newtheorem{proposition}[subsection]{Proposition}
\newtheorem{lemma}[subsection]{Lemma}

\theoremstyle{definition}
\newtheorem{definition}[subsection]{Definition}
\newtheorem{example}[subsection]{Example}
\newtheorem{exercise}[subsection]{Exercise}
\newtheorem{situation}[subsection]{Situation}

\theoremstyle{remark}
\newtheorem{remark}[subsection]{Remark}
\newtheorem{remarks}[subsection]{Remarks}

\numberwithin{equation}{subsection}

% Macros
%
\def\lim{\mathop{\mathrm{lim}}\nolimits}
\def\colim{\mathop{\mathrm{colim}}\nolimits}
\def\Spec{\mathop{\mathrm{Spec}}}
\def\Hom{\mathop{\mathrm{Hom}}\nolimits}
\def\Ext{\mathop{\mathrm{Ext}}\nolimits}
\def\SheafHom{\mathop{\mathcal{H}\!\mathit{om}}\nolimits}
\def\SheafExt{\mathop{\mathcal{E}\!\mathit{xt}}\nolimits}
\def\Sch{\mathit{Sch}}
\def\Mor{\mathop{\mathrm{Mor}}\nolimits}
\def\Ob{\mathop{\mathrm{Ob}}\nolimits}
\def\Sh{\mathop{\mathit{Sh}}\nolimits}
\def\NL{\mathop{N\!L}\nolimits}
\def\CH{\mathop{\mathrm{CH}}\nolimits}
\def\proetale{{pro\text{-}\acute{e}tale}}
\def\etale{{\acute{e}tale}}
\def\QCoh{\mathit{QCoh}}
\def\Ker{\mathop{\mathrm{Ker}}}
\def\Im{\mathop{\mathrm{Im}}}
\def\Coker{\mathop{\mathrm{Coker}}}
\def\Coim{\mathop{\mathrm{Coim}}}

% Boxtimes
%
\DeclareMathSymbol{\boxtimes}{\mathbin}{AMSa}{"02}

%
% Macros for moduli stacks/spaces
%
\def\QCohstack{\mathcal{QC}\!\mathit{oh}}
\def\Cohstack{\mathcal{C}\!\mathit{oh}}
\def\Spacesstack{\mathcal{S}\!\mathit{paces}}
\def\Quotfunctor{\mathrm{Quot}}
\def\Hilbfunctor{\mathrm{Hilb}}
\def\Curvesstack{\mathcal{C}\!\mathit{urves}}
\def\Polarizedstack{\mathcal{P}\!\mathit{olarized}}
\def\Complexesstack{\mathcal{C}\!\mathit{omplexes}}
% \Pic is the operator that assigns to X its picard group, usage \Pic(X)
% \Picardstack_{X/B} denotes the Picard stack of X over B
% \Picardfunctor_{X/B} denotes the Picard functor of X over B
\def\Pic{\mathop{\mathrm{Pic}}\nolimits}
\def\Picardstack{\mathcal{P}\!\mathit{ic}}
\def\Picardfunctor{\mathrm{Pic}}
\def\Deformationcategory{\mathcal{D}\!\mathit{ef}}
\usepackage{cleveref}
\crefname{originallemma}{original lemma}{original lemmas}
\Crefname{originallemma}{Original lemma}{Original lemmas}
\makeatletter
\newlabel{perfect-lemma-affine-compare-bounded}{{06Z0}{source}{}{sourcepointer-1}{}}
\newlabel{perfect-lemma-affine-compare-bounded@cref}{{[originallemma][0][]06Z0}{[1][0][]source}{}{}{}}
\makeatother
\crefname{originaldefinition}{original definition}{original definitions}
\Crefname{originaldefinition}{Original definition}{Original definitions}
\makeatletter
\newlabel{dualizing-definition-relative-dualizing-complex}{{0E2C}{source}{}{sourcepointer-2}{}}
\newlabel{dualizing-definition-relative-dualizing-complex@cref}{{[originaldefinition][0][]0E2C}{[1][0][]source}{}{}{}}
\makeatother
\crefname{originallemma}{original lemma}{original lemmas}
\Crefname{originallemma}{Original lemma}{Original lemmas}
\makeatletter
\newlabel{perfect-lemma-pseudo-coherent}{{08E5}{source}{}{sourcepointer-3}{}}
\newlabel{perfect-lemma-pseudo-coherent@cref}{{[originallemma][0][]08E5}{[1][0][]source}{}{}{}}
\makeatother
\crefname{originallemma}{original lemma}{original lemmas}
\Crefname{originallemma}{Original lemma}{Original lemmas}
\makeatletter
\newlabel{perfect-lemma-quasi-coherence-pushforward}{{0DJK}{source}{}{sourcepointer-4}{}}
\newlabel{perfect-lemma-quasi-coherence-pushforward@cref}{{[originallemma][0][]0DJK}{[1][0][]source}{}{}{}}
\makeatother
\crefname{originallemma}{original lemma}{original lemmas}
\Crefname{originallemma}{Original lemma}{Original lemmas}
\makeatletter
\newlabel{perfect-lemma-quasi-coherence-pullback}{{08DW}{source}{}{sourcepointer-5}{}}
\newlabel{perfect-lemma-quasi-coherence-pullback@cref}{{[originallemma][0][]08DW}{[1][0][]source}{}{}{}}
\makeatother
\crefname{originallemma}{original lemma}{original lemmas}
\Crefname{originallemma}{Original lemma}{Original lemmas}
\makeatletter
\newlabel{perfect-lemma-quasi-coherence-internal-hom}{{0A6H}{source}{}{sourcepointer-6}{}}
\newlabel{perfect-lemma-quasi-coherence-internal-hom@cref}{{[originallemma][0][]0A6H}{[1][0][]source}{}{}{}}
\makeatother
\crefname{originallemma}{original lemma}{original lemmas}
\Crefname{originallemma}{Original lemma}{Original lemmas}
\makeatletter
\newlabel{more-algebra-lemma-more-relative-pseudo-coherent-is-moot}{{0E1T}{source}{}{sourcepointer-7}{}}
\newlabel{more-algebra-lemma-more-relative-pseudo-coherent-is-moot@cref}{{[originallemma][0][]0E1T}{[1][0][]source}{}{}{}}
\makeatother
\title{A relative dualizing complex restricts on an affine chart to the algebraic relative dualizing complex}
\author{The Stacks Project Authors}
\date{stacks-project:duality.tex:local-snapshot-20261012; GNU FDL 1.2 or later}
\begin{document}
\maketitle
\noindent\textit{Editorial note.} All delivered mathematical text is the author's own, copied from the pinned Stacks Project chapter file duality.tex; no statement or proof step was rewritten, reordered or removed.\par\medskip
\noindent\textit{Editorial note.} Preamble provenance: the pinned chapter file opens with an input of the shared preamble and carries no preamble text of its own, so the project's single Stacks-wide preamble staged at sources/stacks-duality/preamble.tex is excerpted verbatim under the recipe fixed by LANE\_BRIEF 22.1, as three spans, lines 8-13 (verbatim package and the reference/slogan/history environments), lines 14-18 (Xy-pic setup) and lines 154-217 (the native theorem declarations and mathematical macros). Nothing was added, removed or redefined; the wrapper class and the amsmath/amssymb/amsthm/graphicx base packages are supplied by the assembler.\par\medskip
\noindent\textit{Editorial note.} The counted statement and proof contain no cross reference, no citation and no included file, so the delivered body is self contained.\par\medskip
\noindent\textit{Editorial note.} Canonical location of the counted proposition: Stacks tag 0E2U (sources/stacks-duality/pinned-original-source.tex line 7271).\par\medskip
\noindent\textit{Allowed standard prerequisites (editorial).}\begin{itemize}
\item The affine comparison between complexes of modules and complexes of sheaves. Derived Categories of Schemes, Lemma labelled perfect-lemma-affine-compare-bounded, cited in the statement. Identifies the restriction of K with an object of D(A).
\item Compatibility of pushforward, pullback and internal Hom with the affine comparison. Derived Categories of Schemes, Lemmas labelled perfect-lemma-quasi-coherence-pushforward, perfect-lemma-quasi-coherence-pullback and perfect-lemma-quasi-coherence-internal-hom, cited at proof lines 7326-7336. Translate the sheaf conditions into algebraic ones.
\item Pseudo-coherence of A over A otimes\_R A for the diagonal. More on Algebra, Lemma labelled more-algebra-lemma-more-relative-pseudo-coherent-is-moot, cited at proof lines 7340-7344. Supplies the pseudo-coherence needed for the internal Hom comparison.
\item The algebraic definition of a relative dualizing complex. Dualizing Complexes, Definition labelled dualizing-definition-relative-dualizing-complex, cited in the statement and at proof lines 7344-7348. Fixes the conditions recovered at the end.
\end{itemize}\medskip
\ifdefined\hypertarget\hypertarget{sourcepointer-1}{}\fi
\noindent\textit{Source reference.} Original lemma 06Z0: sources/stacks-chow/tags.txt, source lines 8904--8904. Cross-chapter standard prerequisite located in the project tag list staged in the local snapshot; invoked at proof line None of the counted proof. It is retained only as a navigation pointer; every load-bearing step is inside the delivered span.\par\medskip
\ifdefined\hypertarget\hypertarget{sourcepointer-2}{}\fi
\noindent\textit{Source reference.} Original definition 0E2C: sources/stacks-chow/tags.txt, source lines 17602--17602. Cross-chapter standard prerequisite located in the project tag list staged in the local snapshot; invoked at proof line 7314 of the counted proof. It is retained only as a navigation pointer; every load-bearing step is inside the delivered span.\par\medskip
\ifdefined\hypertarget\hypertarget{sourcepointer-3}{}\fi
\noindent\textit{Source reference.} Original lemma 08E5: sources/stacks-chow/tags.txt, source lines 10661--10661. Cross-chapter standard prerequisite located in the project tag list staged in the local snapshot; invoked at proof line 7295 of the counted proof. It is retained only as a navigation pointer; every load-bearing step is inside the delivered span.\par\medskip
\ifdefined\hypertarget\hypertarget{sourcepointer-4}{}\fi
\noindent\textit{Source reference.} Original lemma 0DJK: sources/stacks-chow/tags.txt, source lines 16980--16980. Cross-chapter standard prerequisite located in the project tag list staged in the local snapshot; invoked at proof line 7301 of the counted proof. It is retained only as a navigation pointer; every load-bearing step is inside the delivered span.\par\medskip
\ifdefined\hypertarget\hypertarget{sourcepointer-5}{}\fi
\noindent\textit{Source reference.} Original lemma 08DW: sources/stacks-chow/tags.txt, source lines 10652--10652. Cross-chapter standard prerequisite located in the project tag list staged in the local snapshot; invoked at proof line 7302 of the counted proof. It is retained only as a navigation pointer; every load-bearing step is inside the delivered span.\par\medskip
\ifdefined\hypertarget\hypertarget{sourcepointer-6}{}\fi
\noindent\textit{Source reference.} Original lemma 0A6H: sources/stacks-chow/tags.txt, source lines 12845--12845. Cross-chapter standard prerequisite located in the project tag list staged in the local snapshot; invoked at proof line 7303 of the counted proof. It is retained only as a navigation pointer; every load-bearing step is inside the delivered span.\par\medskip
\ifdefined\hypertarget\hypertarget{sourcepointer-7}{}\fi
\noindent\textit{Source reference.} Original lemma 0E1T: sources/stacks-chow/tags.txt, source lines 17583--17583. Cross-chapter standard prerequisite located in the project tag list staged in the local snapshot; invoked at proof line 7310 of the counted proof. It is retained only as a navigation pointer; every load-bearing step is inside the delivered span.\par\medskip
\begin{lemma}
\label{lemma-relative-dualizing-complex-algebra}
Let $X \to S$ be a morphism of schemes which is flat and
locally of finite presentation. Let $(K, \xi)$ be a relative
dualizing complex. Then for any commutative diagram
$$
\xymatrix{
\Spec(A) \ar[d] \ar[r] & X \ar[d] \\
\Spec(R) \ar[r] & S
}
$$
whose horizontal arrows are open immersions, the
restriction of $K$ to $\Spec(A)$ corresponds via
Derived Categories of Schemes, Lemma
\ref{perfect-lemma-affine-compare-bounded}
to a relative dualizing complex for $R \to A$
in the sense of Dualizing Complexes, Definition
\ref{dualizing-definition-relative-dualizing-complex}.
\end{lemma}

\begin{proof}
Since formation of $R\SheafHom$ commutes with restrictions to
opens we may as well assume $X = \Spec(A)$ and $S = \Spec(R)$.
Observe that relatively perfect objects of $D(\mathcal{O}_X)$
are pseudo-coherent and hence are in $D_\QCoh(\mathcal{O}_X)$
(Derived Categories of Schemes, Lemma \ref{perfect-lemma-pseudo-coherent}).
Thus the statement makes sense.
Observe that taking $\Delta_*$, $L\text{pr}_1^*$, and
$R\SheafHom$ is compatible with what happens on the algebraic
side by
Derived Categories of Schemes,
Lemmas \ref{perfect-lemma-quasi-coherence-pushforward},
\ref{perfect-lemma-quasi-coherence-pullback},
\ref{perfect-lemma-quasi-coherence-internal-hom}.
For the last one we observe that $L\text{pr}_1^*K$
is $S$-perfect (hence bounded below) and that $\Delta_*\mathcal{O}_X$
is a pseudo-coherent object of $D(\mathcal{O}_W)$;
translated into algebra this means that $A$ is pseudo-coherent
as an $A \otimes_R A$-module which follows from
More on Algebra, Lemma
\ref{more-algebra-lemma-more-relative-pseudo-coherent-is-moot}
applied to $R \to A \otimes_R A \to A$.
Thus we recover exactly the conditions in
Dualizing Complexes, Definition
\ref{dualizing-definition-relative-dualizing-complex}.
\end{proof}

\end{document}
